We generated 32 two-return and 64 single-return acquisitions for each of nine conditions. The detector always retains the nominal bright-hidden model. True detector recovery changes by $\pm5\%$, electronics recovery by $\pm2\%$, or measured timestamps receive Gaussian noise with standard deviation 0.1 or 0.5 bin. Two further conditions add avalanche-triggered afterpulsing with probabilities 1\% and 5\%. Each accepted avalanche can schedule a descendant after detector recovery plus an exponential waiting time of mean five bins. The electronics filters all avalanches in the same way as in the nominal simulator.

Table~\ref{tab:robust} reports every condition. A measured gap below the nominal support makes the full likelihood zero for every fitted rate, preventing both full-score testing and the common full-null calibration. The table reports their number in the final column. The own-null censored pipeline also receives a separate evaluation with its own fitted null generator. Its known-active criterion remains computable for short gaps; only timestamps outside the acquisition interval prevent its application. Correct-localization rates use all 32 positive acquisitions, so an unavailable record contributes no success. False-alarm rates use applicable nulls, with availability reported separately.

The main sensitivity is the support boundary. When true electronics recovery is 2\% shorter than assumed, the fixed model assigns zero density to all 96 acquisitions. Jitter of 0.1 and 0.5 bin produces a gap below the assumed recovery in 19 and 74 acquisitions, respectively. At 0.5-bin jitter, the full method correctly localizes 4/32 positives, compared with 6/32 for own-null censored calibration. The latter remains applicable to 95 of the 96 acquisitions. Section~\ref{sec:measured} estimates a variable-recovery law directly from measured timestamps.

With true detector recovery 5\% below and above nominal, the full method gives 8/32 and 9/32 correct localizations. With afterpulse probabilities of 1\% and 5\%, it gives 8/32 and 9/32, compared with 2/32 and 3/32 for own-null censored calibration. Its null alarms are 2/64 in each afterpulse condition, versus 1/64 for that control. Compared with the full method's nominal count of 0/64, afterpulsing raises its null alarms to 2/64 at both tested probabilities.

\begin{table*}[t]\centering
\caption{Bright-hidden model-mismatch results at $\alpha=0.01$. Correct localization uses all 32 positive acquisitions, including unavailable records. Wrong locations are counts on those positives. False-alarm fractions use applicable null records. The final column counts positive/null acquisitions for which no test is computed. Common calibration requires the full likelihood; own-null censored calibration uses its own null fit. Dashes indicate no applicable null records.}
\label{tab:robust}
\begin{tabular}{llrrrr}\toprule
True acquisition & Method / calibration & Correct /32 & Wrong & False / applicable & No test (pos; null)\\\midrule
Nominal & Full bank / common & 13/32 & 0 & 0/64 & 0/32; 0/64\\
Nominal & Censored / common & 8/32 & 0 & 2/64 & 0/32; 0/64\\
Nominal & Censored / own-null & 9/32 & 0 & 1/64 & 0/32; 0/64\\
$T_d$: $-5\%$ & Full bank / common & 8/32 & 1 & 1/64 & 0/32; 0/64\\
$T_d$: $-5\%$ & Censored / common & 2/32 & 1 & 1/64 & 0/32; 0/64\\
$T_d$: $-5\%$ & Censored / own-null & 1/32 & 1 & 1/64 & 0/32; 0/64\\
$T_d$: $+5\%$ & Full bank / common & 9/32 & 0 & 1/64 & 0/32; 0/64\\
$T_d$: $+5\%$ & Censored / common & 7/32 & 0 & 1/64 & 0/32; 0/64\\
$T_d$: $+5\%$ & Censored / own-null & 6/32 & 0 & 0/64 & 0/32; 0/64\\
$T_e$: $-2\%$ & Full bank / common & 0/32 & 0 & --- & 32/32; 64/64\\
$T_e$: $-2\%$ & Censored / common & 0/32 & 0 & --- & 32/32; 64/64\\
$T_e$: $-2\%$ & Censored / own-null & 7/32 & 1 & 0/64 & 0/32; 0/64\\
$T_e$: $+2\%$ & Full bank / common & 10/32 & 1 & 0/64 & 0/32; 0/64\\
$T_e$: $+2\%$ & Censored / common & 5/32 & 1 & 1/64 & 0/32; 0/64\\
$T_e$: $+2\%$ & Censored / own-null & 4/32 & 1 & 0/64 & 0/32; 0/64\\
Jitter: 0.1 bin & Full bank / common & 11/32 & 1 & 0/53 & 8/32; 11/64\\
Jitter: 0.1 bin & Censored / common & 5/32 & 1 & 1/53 & 8/32; 11/64\\
Jitter: 0.1 bin & Censored / own-null & 6/32 & 2 & 1/64 & 0/32; 0/64\\
Jitter: 0.5 bin & Full bank / common & 4/32 & 0 & 0/15 & 25/32; 49/64\\
Jitter: 0.5 bin & Censored / common & 1/32 & 0 & 0/15 & 25/32; 49/64\\
Jitter: 0.5 bin & Censored / own-null & 6/32 & 1 & 0/64 & 1/32; 0/64\\
Afterpulse: 1\% & Full bank / common & 8/32 & 0 & 2/64 & 0/32; 0/64\\
Afterpulse: 1\% & Censored / common & 2/32 & 1 & 1/64 & 0/32; 0/64\\
Afterpulse: 1\% & Censored / own-null & 2/32 & 1 & 1/64 & 0/32; 0/64\\
Afterpulse: 5\% & Full bank / common & 9/32 & 0 & 2/64 & 0/32; 0/64\\
Afterpulse: 5\% & Censored / common & 3/32 & 1 & 0/64 & 0/32; 0/64\\
Afterpulse: 5\% & Censored / own-null & 3/32 & 1 & 1/64 & 0/32; 0/64\\
\bottomrule\end{tabular}\end{table*}

